\cleardoublepage
\vspace{2cm}
\begin{center}
{\fontsize{14}{16.8}\selectfont\textbf{DAFTAR SIMBOL}}\\[1em]
\end{center}
\label{chap:simbol}
\addcontentsline{toc}{chapter}{DAFTAR SIMBOL}
\setcounter{chapter}{0}
\setcounter{section}{0}
\vspace{2em}




\begin{longtable}{@{} l >{\raggedright\arraybackslash}p{6.5cm} c @{}}
\multicolumn{1}{l}{\textbf{Notasi}} & \multicolumn{1}{l}{\textbf{Deskripsi}} & \multicolumn{1}{l}{\parbox[c]{3.5cm}{\centering\textbf{Pemakaian pertama kali}}} \\
\midrule

% Konsumsi Energi dan Baterai
$EC(v)$ & Fungsi tingkat konsumsi energi (\textit{Energy Consumption}) dalam satuan kWh/km & \pageref{eq:ec-mamarikas} \\
$v$ & Kecepatan operasional bus dalam satuan km/jam & \pageref{eq:ec-mamarikas} \\
$a, b, c, d, e, f, g$ & Koefisien kalibrasi kurva energi spesifik bus pada model Mamarikas & \pageref{eq:ec-mamarikas} \\
$SoC_i$ & \textit{State of Charge} (tingkat daya baterai tersisa) pada saat tiba di perhentian $i$ & \pageref{eq:soc-update} \\
$E_{bat}$ & Total kapasitas maksimum baterai bus dalam satuan kWh & \pageref{eq:soc-update} \\
$d_i$ & Jarak tempuh (metrik spasial) pada segmen perjalanan ke-$i$ dalam satuan km & \pageref{eq:soc-update} \\
$v_i$ & Kecepatan rata-rata spesifik pada segmen ke-$i$ dalam satuan km/jam & \pageref{eq:soc-update} \\[1ex]

% Variasi dan Randomisasi (Monte Carlo)
$t_{\text{base}}$ & Waktu tempuh dasar ideal sebuah segmen perjalanan tanpa efek variasi & \pageref{eq:lognormal-time} \\
$t_{\text{expected}}$ & Ekspektasi waktu tempuh harian setelah dikenakan pengali kemacetan & \pageref{eq:lognormal-time} \\
$t_{\text{aktual}}$ & Waktu tempuh aktual pada satu realisasi simulasi (\textit{random variate}) & \pageref{eq:lognormal-time} \\
$M_{\text{daily}}$ & \textit{Modifier} (faktor pengali) tingkat kemacetan harian sistemik & \pageref{eq:lognormal-time} \\
$CV_{\text{seg}}$ & \textit{Coefficient of Variation} untuk kecepatan level segmen & \pageref{eq:lognormal-time} \\
$\sigma$ & Parameter deviasi standar dari distribusi \textit{log-normal} & \pageref{eq:lognormal-time} \\
$\mu$ & Parameter rata-rata (\textit{mean}) dari distribusi \textit{log-normal} & \pageref{eq:lognormal-time} \\
$N$ & Total jumlah iterasi atau jumlah hari pengujian dalam \textit{Monte Carlo} & \pageref{alg:monte-carlo} \\[1ex]

% Headway dan Layanan
$\bar{h}_r$ & Rata-rata \textit{headway} (waktu jeda antar-keberangkatan) untuk koridor/rute $r$ & \pageref{eq:headway-dev} \\
$\Delta h_r$ & Deviasi \textit{headway} rata-rata rute $r$ terhadap patokan standar Dishub & \pageref{eq:headway-dev} \\
$h_{\text{target},r}$ & Target \textit{headway} resmi yang ditetapkan oleh Dishub DIY untuk rute $r$ & \pageref{eq:headway-dev} \\
$t_{i,r}$ & Waktu riil keberangkatan rit perjalanan ke-$i$ pada rute $r$ & \pageref{eq:headway-dev} \\
$N_r$ & Total jumlah keberangkatan rit yang berhasil diselesaikan pada rute $r$ & \pageref{eq:headway-dev} \\[1ex]

% Simulasi, SPKLU, dan Rute Lumpuh
$B_r$ & Himpunan armada bus yang ditugaskan beroperasi pada rute $r$ & \pageref{eq:degraded-route} \\
$SoC_b$ & Tingkat daya baterai aktual dari armada bus spesifik $b$ & \pageref{eq:degraded-route} \\
$w_{q,b}$ & Durasi tunggu (waktu antrean di SPKLU) yang dialami bus $b$ & \pageref{eq:degraded-route} \\
$t_{\text{cas}}$ & Waktu simulasi riil (menit) saat bus $b$ mulai dialiri listrik di SPKLU & \pageref{eq:degraded-route} \\
$t_{\text{akhir\_shift}}$ & Batas waktu akhir operasional layanan reguler harian (20.30) & \pageref{eq:degraded-route} \\
$Q$ & Antrean peristiwa (\textit{Priority Queue}) pada mesin \textit{Discrete-Event Simulation} & \pageref{alg:des-loop} \\

\end{longtable}
